% มคอ.5 — 4095101 ภาคเรียนที่ 1/2568 (รุ่นแรก)
\documentclass{mko5}
\begin{document}
\mkofivetitle

\coursecode{4095101}
\coursename{ปัญญาประดิษฐ์และการเรียนรู้ของเครื่อง}
\curriculumnote{หลักสูตรวิศวกรรมศาสตรมหาบัณฑิต สาขาวิชาวิศวกรรมคอมพิวเตอร์และปัญญาประดิษฐ์ พ.ศ. 2568}
\coursegroup{วิชาเฉพาะด้านบังคับ}
\coursecredits{3 (2-2-5)}
\studyplan{แผน 2 --- วิชาชีพ (สารนิพนธ์)}
\semester{ภาคเรียนที่ 1 / ปีการศึกษา 2568}
\prerequisite{ไม่มี}
\corequisite{ไม่มี}
\studyplace{มหาวิทยาลัยราชภัฏอุตรดิตถ์ --- คณะเทคโนโลยีอุตสาหกรรม}
\registeredstudents{4 คน}
\completedstudents{4 คน}
\instructor{ผศ.ดร.พิศิษฐ์ นาคใจ และ ผศ.ดร.พัชรี มณีรัตน์}
\coinstructors{---}

\mkoheaderinfo

\begin{teachingactualtable}
  \actualrow{CLO1}{บรรยายทฤษฎี AI/ML ยกตัวอย่างกรณีศึกษา สาธิต อภิปรายกลุ่ม แบบทดสอบและแบบฝึกหัด}
  \actualrow{CLO2}{สาธิตเครื่องมือวิเคราะห์ข้อมูล Workshop โครงงานย่อย การนำเสนอผล}
  \actualrow{CLO3}{สาธิตพัฒนาโมเดล Deep Learning ทดสอบและปรับปรุงโมเดล รายงานผลการทดสอบ}
  \actualrow{CLO4}{อภิปรายจริยธรรม AI บทความสะท้อนคิด การมีส่วนร่วมในชั้นเรียน}
  \actualrow{CLO5}{พัฒนาระบบ AI ทดสอบระบบ นำเสนอผลงาน โครงงานปลายภาค}
\end{teachingactualtable}

\begin{assessmentusedtable}
  \assessusedrow{CLO1}{การสอบข้อเขียน การนำเสนอหน้าชั้น การมีส่วนร่วมในชั้นเรียน}
  \assessusedrow{CLO2}{การสอบปฏิบัติ รายงานผลการวิเคราะห์ การนำเสนอโครงงาน}
  \assessusedrow{CLO3}{โครงงาน Deep Learning รายงานผลการทดสอบ การนำเสนอผลงาน}
  \assessusedrow{CLO4}{การมีส่วนร่วมในการอภิปราย บทความสะท้อนคิด}
  \assessusedrow{CLO5}{โครงงานปลายภาค การสาธิตระบบ รายงานโครงงาน --- Rubric 4 ระดับ}
\end{assessmentusedtable}

\mkosectionthree

\closummaryblock{CLO1}{อธิบาย หลักการพื้นฐานของปัญญาประดิษฐ์และการเรียนรู้ของเครื่อง ทั้งแบบมีผู้สอนและไม่มีผู้สอนได้อย่างถูกต้อง}{80}
\achievementpass{4}{100}{ผ่านเกณฑ์ระดับ 3 ขึ้นไปทุกคน}
\achievementfail{0}{0}{ไม่มี}
\closummaryend

\closummaryblock{CLO2}{ประยุกต์ใช้ เทคนิคการแก้ปัญหาด้วยการค้นหาและการวิเคราะห์ข้อมูล ในการพัฒนาแบบจำลองการเรียนรู้ของเครื่องได้}{80}
\achievementpass{4}{100}{โครงงานและสอบปฏิบัติผ่าน}
\achievementfail{0}{0}{ไม่มี}
\closummaryend

\closummaryblock{CLO3}{วิเคราะห์ ประสิทธิภาพของแบบจำลองโครงข่ายประสาทเทียมและการเรียนรู้เชิงลึก เพื่อเลือกใช้และปรับปรุงให้เหมาะสมกับปัญหาที่กำหนดได้}{80}
\achievementpass{4}{100}{โครงงาน DL ผ่านเกณฑ์}
\achievementfail{0}{0}{ไม่มี}
\closummaryend

\closummaryblock{CLO4}{แสดงออก ถึงความตระหนักในจริยธรรมและความรับผิดชอบต่อสังคม ในการพัฒนาและประยุกต์ใช้ระบบปัญญาประดิษฐ์}{80}
\achievementpass{4}{100}{บทสะท้อนคิดและการอภิปรายผ่าน}
\achievementfail{0}{0}{ไม่มี}
\closummaryend

\closummaryblock{CLO5}{พัฒนา ระบบปัญญาประดิษฐ์ โดยใช้เทคนิคการเรียนรู้ของเครื่องที่เหมาะสมในการแก้ปัญหาจริงได้อย่างมีประสิทธิภาพ}{80}
\achievementpass{4}{100}{โครงงานปลายภาคและการสาธิตระบบผ่าน}
\achievementfail{0}{0}{ไม่มี}
\closummaryend

\mkosectionfour
\begin{mkotextblock}
การเข้าถึง GPU สำหรับการเทรนโมเดลขนาดใหญ่ยังจำกัด;
นักศึกษาบางรายต้องการเวลาเพิ่มในการทำความเข้าใจทฤษฎีเชิงคณิตศาสตร์เบื้องต้น
\end{mkotextblock}

\mkosectionfive
\begin{mkotextblock}
เพิ่มชุดแบบฝึกหัดก่อนเรียน (pre-test) เพื่อจัดกลุ่มช่วยเหลือ;
จัดทำตัวอย่างโค้ดและ Notebook มาตรฐานสำหรับแต่ละหัวข้อ;
เชื่อมโยงกับวิชา 4095102 ในลำดับถัดไป
\end{mkotextblock}

\mkosectionsix
\begin{mkotextblock}
ขอสนับสนุนการเข้าถึง GPU Server สำหรับรายวิชา AI/ML อย่างต่อเนื่อง
\end{mkotextblock}

\mkosectionseven
\mkosignaturepair{ผศ.ดร.พิศิษฐ์ นาคใจ และ ผศ.ดร.พัชรี มณีรัตน์}{30 เมษายน 2569}{ผศ.ดร.พิทักษ์ คล้ายชม}{30 เมษายน 2569}
\end{document}
